% ============================================================
%  chapter_feynpy.tex  ---  User Manual for the Feynman-rule library

% ============================================================

% ==============================================================
\chapter{FeynPy: User Manual}
\label{ch:feynpy}
% ==============================================================

The purpose of this chapter is to serve as a guide to creating a model in \textsc{FeynPy}. Theoretical explanations of the elements that make up the model are provided in the previous chapters, while an overview of how the algorithm works is given in Chapter~\ref{ch:internals}. Here, we focus on the user interface. If we take the following Lagrangian as our starting point:

\begin{equation}
\mathcal{L}_{\mathrm{QCD}}
=
i\,\bar{q}\gamma^\mu D_\mu q
-\frac{1}{4}F_{\mu\nu}^{a}F^{a\,\mu\nu}
-\frac{1}{2\xi}
\left(\partial^\mu A_\mu^{a}\right)^2
-\bar{c}^{a}\,\partial^\mu(D_\mu c)^{a},
\end{equation}
we note that the elements composing this Lagrangian are, first and foremost,
fermionic fields \(q_{i,c,f}\) carrying Dirac, colour, and flavour indices. Furthermore, we have the field-strength tensor $F^a_{\mu\nu}$, which must be expanded as in Equation~\ref{eq:fs-general}, thus giving rise to two point, three point and four point gluon terms. There is also the covariant derivative, whose
expansion introduces the interaction between the fermionic and gauge fields.
Finally, the Lagrangian contains the gauge-fixing and ghost terms. These
expansions are determined by the symmetry group associated with the
Lagrangian. Our aim will therefore be to define the basic ingredients of the model: indices, fields, gauge representations and gauge groups, and then assemble them into a complete \texttt{Model} object. Once the model is defined, the toolkit can be used to manipulate the Lagrangian, extract Feynman rules and perform further symbolic operations.

The chapter is organised in three movements.
\Cref{sec:feynpy-imports,sec:feynpy-indices,sec:feynpy-fields,sec:feynpy-gauge,sec:feynpy-flavor,sec:feynpy-parameters,sec:feynpy-model,sec:feynpy-operators}
build a model from its smallest pieces up to a complete \texttt{Model}.
\Cref{sec:feynpy-feynman,sec:feynpy-discovery,sec:feynpy-qcd-example,sec:feynpy-flavor-expand}
extract vertices from it, ending with the QCD Lagrangian above carried
through to its printed rules. The remaining sections
(\Cref{sec:feynpy-transformations} onwards) cover what can be done to a
compiled Lagrangian besides reading vertices off it: rewriting it into a
physical basis, exporting it to \textsc{Symbolica}, and acting on it with
operators such as an infinitesimal gauge variation or the BRST differential.
\Cref{sec:feynpy-models} closes with the three complete models that ship with
the toolkit.




% --------------------------------------------------------------
\section{Core Imports}
\label{sec:feynpy-imports}
% --------------------------------------------------------------

\begin{minted}{python}
from symbolica import S, Expression
from symbolic.vertex_engine import I
from symbolic.spenso_structures import (
    gauge_generator, structure_constant,
    weak_gauge_generator, weak_structure_constant,
)
from feynpy import (
    Field, GhostField, GaugeGroup, GaugeRepresentation, IndexType,
    Model, Parameter, FieldTransformation,
    # declarative Lagrangian factors
    DC, FS, PartialD, Gamma, Gamma5, Metric, T, StructureConstant,
    GaugeFixing, GhostLagrangian, ProjM, ProjP, rotation,
    # pre-built index constants
    SPINOR_INDEX, LORENTZ_INDEX,
    COLOR_FUND_INDEX, COLOR_ADJ_INDEX,
    WEAK_FUND_INDEX, WEAK_ADJ_INDEX,
    # flavour
    flavor_index,
)
\end{minted}


\noindent Let's start our journey with a quick overview of a minimal complete example, the scalar $\phi^4$ theory:

\begin{minted}{python}
Phi   = Field("Phi", spin=0, self_conjugate=True, symbol=S("phi"))
lam   = S("lam")
model = Model(lam * Phi * Phi * Phi * Phi)
L     = model.lagrangian()
V     = L.feynman_rule(Phi, Phi, Phi, Phi)   # -> 24i*lam
\end{minted}

\noindent Four lines: declare a field, write a Lagrangian, compile it, ask for
a vertex. Everything that follows in this chapter is an elaboration of those
four lines.

% --------------------------------------------------------------
\section{Index Types}
\label{sec:feynpy-indices}
% --------------------------------------------------------------

Many fields carry indices that define how they transform under a symmetry. For example, a quark $q_{i,c,f}$ carries three indices:
\begin{enumerate}
    \item a spinor index $i$,
    \item a colour fundamental index $c$,
    \item a flavour index $f$.
\end{enumerate}
In the toolkit, indices are represented by objects of type \texttt{IndexType}.
Internally, each \texttt{IndexType} stores the index name, pairs a \textsc{Spenso} \py{Representation} with a string \emph{kind} that is used to match indices through the compilation pipeline, and the prefix used when dummy indices are generated automatically. Six pre-built constants cover the standard cases and should be used directly whenever possible.




\begin{center}
\begin{tabular}{@{}P{3.8cm}ll@{}}
\toprule
\normalfont\textbf{Constant} & \textbf{Representation} & \textbf{kind string} \\
\midrule
\texttt{SPINOR\_INDEX}     & \texttt{bis(4)}  & \texttt{"spinor"}     \\
\texttt{LORENTZ\_INDEX}    & \texttt{mink(4)} & \texttt{"lorentz"}    \\
\texttt{COLOR\_FUND\_INDEX}& \texttt{cof(3)}  & \texttt{"color\_fund"}\\
\texttt{COLOR\_ADJ\_INDEX} & \texttt{coad(8)} & \texttt{"color\_adj"} \\
\texttt{WEAK\_FUND\_INDEX} & \texttt{cof(2)}  & \texttt{"weak\_fund"} \\
\texttt{WEAK\_ADJ\_INDEX}  & \texttt{coad(3)} & \texttt{"weak\_adj"}  \\
\bottomrule
\end{tabular}
\end{center}

\noindent All six can be imported directly from \texttt{feynpy}. The colour adjoint index type, for instance, is defined as:
\begin{minted}{python}
COLOR_ADJ_INDEX = IndexType(
    "ColorAdj", Representation.coad(8), "color_adj", prefix="a"
)
\end{minted}
\noindent Custom index types follow the same pattern, so a new gauge group
with a representation Spenso already knows requires no engine changes:

\begin{minted}{python}
from symbolica.community.spenso import Representation
from feynpy import IndexType

# IndexType(name, representation, kind, prefix=...)
MY_FUND = IndexType("MyFund", Representation.cof(3), "my_fund", prefix="m")
\end{minted}

\noindent Flavor indices are ordinary \py{IndexType} objects with
\py{is_flavor=True} set automatically by the helper:

\begin{minted}{python}
# flavor_index(name, dimension, prefix=...)
Generation = flavor_index("Generation", 3, prefix="f")
\end{minted}

% --------------------------------------------------------------
\section{Field Declaration}
\label{sec:feynpy-fields}
% --------------------------------------------------------------

Once the relevant indices have been defined, we can move on to defining the fields of the model. The structure of the fields are insipired by the particle-class declarations used in FeynRules but expressed by a Python object. A particle field will be characterised by the following properties: name, spin, statistics, index structure, quantum numbers, and, when needed, its mass and conjugation properties.

In \textsc{FeynPy} fields are represented by the \texttt{Field} class, whose properties are listed in the following table.

\begin{center}
\begin{tabular}{@{}P{3.8cm}p{8.5cm}@{}}
\toprule
\normalfont\textbf{Parameter} & \textbf{Description} \\
\midrule
\texttt{name}              & Symbolic name of the field. \\
\texttt{spin}              & \py{0}, \py{Fraction(1,2)}, or \py{1}. \\
\texttt{self\_conjugate}   & \py{True} if the field is equal to its conjugate. \\
\texttt{symbol}            & Symbolica atom for this field. \\
\texttt{conjugate\_symbol} & Required when \py{self_conjugate=False}. \\
\texttt{indices}           & Tuple of \py{IndexType} objects carried by the field. \\
\texttt{quantum\_numbers}  & Dict of charge labels, e.g.\ \mintinline{python}|{"Q": S("e")}|. Required for abelian gauge coupling. \\
\texttt{kind}              & Set to \py{"ghost"} for ghost fields. \\
\texttt{flavor\_index}     & The \py{IndexType} to use for flavor expansion. \\
\texttt{class\_members}    & Tuple of member-name strings, one per generation. \\
\bottomrule
\end{tabular}
\end{center}


\noindent
Note that not every entry is required in every case: a real scalar needs almost nothing, while a quark needs a spinor index, a colour index and a conjugate symbol. The example below declares the four fields of the QCD Lagrangian used as the running example of this chapter, plus one charged scalar to show how an abelian charge is attached.

\begin{minted}{python}
from fractions import Fraction

# Quark - spinor and fundamental colour index
Quark  = Field("q", spin=Fraction(1,2), self_conjugate=False,
               symbol=S("q"), conjugate_symbol=S("qbar"),
               indices=(SPINOR_INDEX, COLOR_FUND_INDEX))

# Gluon - non-abelian gauge boson, Lorentz and adjoint index
Gluon  = Field("G", spin=1, self_conjugate=True,
               symbol=S("G"), indices=(LORENTZ_INDEX, COLOR_ADJ_INDEX))

# Faddeev-Popov ghost of the gluon
GhostG = GhostField("ghG", ghost_of=Gluon, self_conjugate=False,
                    symbol=S("ghG"), conjugate_symbol=S("ghGbar"),
                    indices=(COLOR_ADJ_INDEX,))

# Complex scalar carrying a U(1) charge, for comparison
PhiC   = Field("PhiC", spin=0, self_conjugate=False,
               symbol=S("phiC"), conjugate_symbol=S("phiCdag"),
               quantum_numbers={"Q": S("Qphi")})
\end{minted}

\noindent Two points are worth stressing. A gauge boson is an ordinary \py{Field}; it becomes a gauge boson only by being named as such in a \py{GaugeGroup}. And an abelian coupling is read from \py{quantum_numbers}, so a field with no entry for the relevant charge key simply does not couple to that group.

\begin{tcolorbox}[notebox, title=Note]
The conjugated version of any non-self-conjugate field is written with \texttt{.bar}, for example \py{Psi.bar}, \py{Higgs.bar}.  Use this form both in Lagrangian expressions and in \py{feynman_rule(...)} calls. 
\end{tcolorbox}

% --------------------------------------------------------------
\section{Gauge Groups and Representations}
\label{sec:feynpy-gauge}
% --------------------------------------------------------------
Before defining the full model, more physical information needs to be declared. In particular, the symmetry structure of the model must be specified through the gauge representation and the gauge group.

To define a gauge representation, three elements are required: the index, the generator, and a name of your choice. For fields carrying multiple indices of the same kind, the optional \texttt{slot} argument specifies which index slot is acted on by the generator, while \texttt{slot\_policy="sum"} instructs the toolkit to sum the generator action over all matching slots rather than requiring a unique one. After creating the \texttt{GaugeRepresentation}, one can assemble the \texttt{GaugeGroup}.
\subsection*{GaugeRepresentation parameters}

\begin{center}
\begin{tabular}{@{}P{3.3cm}p{9cm}@{}}
\toprule
\normalfont\textbf{Parameter} & \textbf{Description} \\
\midrule
\texttt{index}            & The \py{IndexType} for this representation. \\
\texttt{generator\_builder} & Function that builds the generator tensor: \py{gauge_generator} for the built-in color $\mathrm{SU}(3)$ fundamental, \py{weak_gauge_generator} for $\mathrm{SU}(2)$ doublets, or a custom builder for other groups. \\
\texttt{name}             & Label string. \\
\texttt{slot}             & For fields with repeated same-kind indices, specifies which slot (0-based) is active. \\
\texttt{slot\_policy}     & \py{"sum"} to sum over all matching slots instead of requiring uniqueness. \\
\bottomrule
\end{tabular}
\end{center}

\subsection*{GaugeGroup parameters}


\begin{center}
\begin{tabular}{@{}P{3.3cm}p{9cm}@{}}
\toprule
\normalfont\textbf{Parameter} & \textbf{Description} \\
\midrule
\texttt{name}               & Identifier string. \\
\texttt{abelian}            & \py{True} for abelian groups. \\
\texttt{coupling}           & Symbolica symbol for the coupling constant. \\
\texttt{gauge\_boson}       & The gauge boson \py{Field}. \\
\texttt{charge}             & For abelian groups: the quantum-number key to read from matter fields (e.g.\ \py{"Q"}). \\
\texttt{structure\_constant}& For non-abelian groups: the structure-constant builder (e.g.\ \py{structure_constant}). \\
\texttt{representations}    & Tuple of \py{GaugeRepresentation} objects for matter fields. \\
\texttt{ghost\_field}       & The ghost \py{Field} --- required for \py{GhostLagrangian(...)}. \\
\bottomrule
\end{tabular}
\end{center}

Having defined the indices and representations, the symmetry group can now be defined. As an example, we illustrate how to define the symmetry groups of the SM: $SU(3)_C \times SU(2)_L \times U(1)_Y$

Note that the abelian group $U(1)_Y$ is specified by its charge, while the non-abelian groups $SU(3)_C$ and $SU(2)_L$ require structure constants and a list of representations.
\begin{minted}{python}
# Representation for the colour fundamental
COLOR_FUND_REP = GaugeRepresentation(index=COLOR_FUND_INDEX,
    generator_builder=gauge_generator, name="fundamental")

# U(1)_Y
U1Y = GaugeGroup(name="U1Y", abelian=True,
                coupling=S("g"), gauge_boson=B, charge="Y")

# SU(3) QCD
SU3C  = GaugeGroup(name="SU3C",  abelian=False, coupling=S("gs"),
                   gauge_boson=Gluon, ghost_field=GhostG,
                   structure_constant=structure_constant,
                   representations=(COLOR_FUND_REP,))

# SU(2) weak -- doublet representation
WEAK_DOUBLET_REP = GaugeRepresentation(index=WEAK_FUND_INDEX,
    generator_builder=weak_gauge_generator, name="doublet")

SU2L = GaugeGroup(name="SU2L", abelian=False, coupling=S("g2"),
                  gauge_boson=W, ghost_field=GhostW,
                  structure_constant=weak_structure_constant,
                  representations=(WEAK_DOUBLET_REP,))
\end{minted}

% --------------------------------------------------------------
\section{Flavor Indices}
\label{sec:feynpy-flavor}
% --------------------------------------------------------------

Let's move on to falvours. The three generations of quarks and leptons are copies of the same field, and writing them out separately would triple the length of every Lagrangian.
Flavor indices allow compact multi-generation expressions that can later be expanded into per-member Feynman rules. A field declares which of its indices is the flavour one, and optionally the names of its class members.

\begin{minted}{python}
# 1.  Create a flavour index of dimension 3
Generation = flavor_index("Generation", 3, prefix="f")

# 2.  Field with a flavour index and explicit class members
lepton = Field("l", spin=Fraction(1,2), self_conjugate=False,
               indices=(Generation, SPINOR_INDEX),
               symbol=S("l"), conjugate_symbol=S("lbar"),
               flavor_index=Generation,
               class_members=("e", "mu", "ta"))
e, mu, ta = lepton.class_members
\end{minted}

\noindent The pattern mirrors the \texttt{ClassMembers} entry of a \textsc{FeynRules} particle class. 

% --------------------------------------------------------------
\section{Parameters}
\label{sec:feynpy-parameters}
% --------------------------------------------------------------

Couplings, masses and mixing matrices are declared as \py{Parameter} objects. A \py{Parameter} behaves like its \textsc{Symbolica} symbol in algebraic expressions, so a bare symbol from \py{S("lam")} works just as well for a plain constant; the class earns its place when the parameter carries indices or has known components.

\begin{center}
\begin{tabular}{@{}P{3.4cm}p{9cm}@{}}
\toprule
\normalfont\textbf{Parameter} & \textbf{Description} \\
\midrule
\texttt{name}             & Name of the parameter. \\
\texttt{symbol}           & Symbolica symbol; defaults to \py{S(name)}. \\
\texttt{indices}          & Tuple of \py{IndexType} objects it carries, e.g.\ two flavour indices for a Yukawa matrix. \\
\texttt{complex\_param}   & \py{True} if the parameter is complex; controls conjugation. \\
\texttt{components}       & Dict mapping index tuples to fixed values; a value of \py{0} removes that component entirely. \\
\texttt{unitary\_partner} & The conjugate matrix, for unitary pairs such as CKM and CKM$^\dagger$. \\
\texttt{value}            & Optional numeric or symbolic value, stored as metadata. \\
\bottomrule
\end{tabular}
\end{center}

\begin{minted}{python}
# Plain coupling constant
lam = Parameter("lam")

# Complex Yukawa matrix on two flavour indices
Yu  = Parameter("Yu", indices=(Generation, Generation), complex_param=True)

# Diagonal matrix: every off-diagonal component is forced to zero
Ydiag = Parameter("Ydiag", indices=(Generation, Generation),
                  components={(i, j): 0
                              for i in range(1, 4)
                              for j in range(1, 4) if i != j})

# Unitary pair, used by rotation(...) in field transformations
CKM    = Parameter("CKM", indices=(Generation, Generation),
                   complex_param=True)
CKMDag = Parameter("CKMDag", indices=(Generation, Generation),
                   complex_param=True, unitary_partner=CKM)
\end{minted}

\noindent An indexed parameter is used by calling it: \py{Yu(f, g)} inside a
Lagrangian term contracts its two flavour indices against the fields that
carry the labels \py{f} and \py{g}. The \py{components} mechanism is what
makes \py{Ydiag} produce no off-diagonal vertices at all.


% --------------------------------------------------------------
\section{Declaring a Model}
\label{sec:feynpy-model}
% --------------------------------------------------------------
Our journey goes to the focal point, we can put everything together to start then calculation and manipulation. After declaring indices, fields, parameters, representations and gauge groups, the full model can be constructed. It is represented as a Python class containing the following information: the name, the gauge groups, the parameters and the Lagrangian.

For simpler models, it is sufficient to write the Lagrangian. For example, a quartic scalar interaction can be defined as
\begin{equation}
\mathcal{L}_{\phi^4} = \lambda \phi^4 .
\end{equation}

If the Lagrangian contains covariant derivatives or field-strength tensor, it will be necessary to specify the gauge group so that the toolkit will know the generator, couplings, gauge bosons and structure constant. For example, the QED fermion Lagrangian is

\begin{equation}
\mathcal{L}_{\mathrm{QED}}
= i \bar{\psi}\gamma^\mu D_\mu\psi ,
\end{equation}

while the complete gauge-fixed QCD Lagrangian, that we wanted to compute at the beginning of the chapter, is

\begin{equation}
\mathcal{L}_{\mathrm{QCD}}
=
i\bar{q}\gamma^\mu D_\mu q
- \frac{1}{4}F_{\mu\nu}^{a}F^{a\mu\nu}
+ \mathcal{L}_{\mathrm{GF}}
+ \mathcal{L}_{\mathrm{ghost}} .
\end{equation}

These models are implemented as follows:


\begin{minted}{python}
# Local model - no gauge_groups needed
m = Model(S("lam") * Phi * Phi * Phi * Phi)

# Gauge-aware model
mu = S("mu")
m  = Model(
    I * Psi.bar * Gamma(mu) * DC(Psi, mu),
    gauge_groups=(U1QED,),
)

# Multi-term model (full gauge-fixed QCD)
mu, nu, a, xi = S("mu"), S("nu"), S("a"), S("xi")
qcd = Model(
    I * Quark.bar * Gamma(mu) * DC(Quark, mu)
    - (Expression.num(1)/Expression.num(4))
      * FS(SU3C, mu, nu, a) * FS(SU3C, mu, nu, a)
    + GaugeFixing(SU3C, xi=xi)
    + GhostLagrangian(SU3C),
    gauge_groups=(SU3C,),
)

# Compile
L = qcd.lagrangian()
\end{minted}

\noindent The four summands of the last model correspond one-to-one to the four terms of the QCD Lagrangian written at the start of this chapter. \Cref{sec:feynpy-qcd-example} carries this example through to its vertices.

\subsection*{Model validation}

\begin{minted}{python}
# Check gauge-sector consistency before extracting rules
report = m.validate()
print(report.ok)      # True if no issues
print(report.issues)  # list of Issue(code, message)
\end{minted}

% --------------------------------------------------------------
\section{Operator Catalog}
\label{sec:feynpy-operators}
% --------------------------------------------------------------

Before going to how to manipulate this lagrangians lets first see common operator that we can find as valid Lagrangian building block inside \py{Model(...)}.

\subsection*{Local monomials}

\begin{minted}{python}
Field * Field * ...                          # plain product
Psi.bar * Psi * Phi                          # Yukawa
Psi.bar * Gamma(mu) * Psi * Photon           # vector current
Psi.bar * Gamma(mu) * Gamma5() * Psi * V    # axial current
\end{minted}

\subsection*{Derivative operators --- \texttt{PartialD}}

\begin{minted}{python}
PartialD(Phi, mu)                            # partial_mu Phi
PartialD(PartialD(Phi, mu), nu)              # nested derivative
Psi.bar * Gamma(mu) * PartialD(Psi, nu) * Phi
\end{minted}

\subsection*{Gauge covariant derivative --- \texttt{DC}}

Convention: $D_\mu = \partial_\mu - igA_\mu$, as fixed in
\Cref{tab:theory-conventions}.

\begin{minted}{python}
mu = S("mu")
# Fermion kinetic term
I * Psi.bar * Gamma(mu) * DC(Psi, mu)

# Complex scalar kinetic term
DC(PhiC.bar, mu) * DC(PhiC, mu)

# Restrict the expansion to one group of a multi-charged field
DC(QL, mu, SU3C)
\end{minted}

\subsection*{Field strength --- \texttt{FS}}

\begin{minted}{python}
# Abelian (no adjoint index)
FS(U1QED, mu, nu)

# Non-abelian (adjoint index required)
FS(SU3C, mu, nu, S("a"))

# Yang-Mills kinetic term
-(Expression.num(1)/Expression.num(4)) \
    * FS(SU3C, mu, nu, S("a")) \
    * FS(SU3C, mu, nu, S("a"))
\end{minted}

\begin{tcolorbox}[warnbox, title={Strict rules for \texttt{FS}}]
A non-abelian \py{FS} requires an explicit adjoint index; an abelian one must
not carry one.  All adjoint indices must be contracted, open colour
indices raise a \py{ValueError} at compile time.
\end{tcolorbox}

\noindent Higher powers and mixed-group:

\begin{minted}{python}
from symbolic.spenso_structures import lorentz_levi_civita

# CP-odd operator built from the dual field strength
theta * lorentz_levi_civita(mu, nu, rho, sig) \
      * FS(U1QED, mu, nu) * FS(U1QED, rho, sig)

# Cubic field-strength operator
c3 * StructureConstant(SU3C, a, b, c) \
   * FS(SU3C, mu, nu, a) * FS(SU3C, nu, rho, b) * FS(SU3C, rho, mu, c)
\end{minted}

\subsection*{Gauge fixing and ghosts helper}

These are helper used to compile the gauge fixing and ghost lagrangian given a gauge group.
\begin{minted}{python}
GaugeFixing(SU3C, xi=S("xi"))   # linear covariant gauge
GhostLagrangian(SU3C)           # Faddeev-Popov
\end{minted}
The user is free to define it manually as a terml.
% --------------------------------------------------------------
\section{Extracting Feynman Rules}
\label{sec:feynpy-feynman}
% --------------------------------------------------------------

We finally have our model describe the Lagrangain and the symmetry. The purpose of this project was to achieve the Feynman Rule given a Lagrangian and now we have everything we need to extract the rules and do much more.

Call \py{feynman_rule} on the \py{Model} or the compiled \py{Lagrangian} object returned by
\py{model.lagrangian()} to obtain the Feynman Rule expressed as a Symbolica expression.

\begin{minted}{python}
L = model.lagrangian()

# Targeted: pass explicit fields to get one vertex
vertex = L.feynman_rule(Psi.bar, Psi, Photon)

# Custom external momenta (default labels: q1, q2, q3, ...)
vertex = L.feynman_rule(Psi.bar, Psi, Photon,
                        momenta=[S("p_in"), S("p_out"), S("k")])
\end{minted}

\subsection*{Key options}

\begin{center}
\begin{tabular}{@{}P{3.3cm}lp{7.5cm}@{}}
\toprule
\normalfont\textbf{Option} & \textbf{Default} & \textbf{Description} \\
\midrule
\texttt{simplify}         & \py{True}  & Run symbolic simplification on the result. \\
\texttt{simplify\_gamma}  & \py{False} & If \py{simplify=True}, also apply Dirac-algebra simplification to gamma chains. \\
\texttt{include\_delta}   & \py{False} & Include the overall momentum-conservation $\delta$ factor. \\
\texttt{strip\_externals} & \py{True}  & Strip external wavefunction factors. \\
\texttt{flavor\_expand}   & \py{False} & \py{True}, one flavor index, or an iterable of flavor indices to expand flavor-class members. \\
\texttt{key\_format}      & \py{"names"} & Result-dict keys: \py{"names"} (string tuples) or \py{"fields"} (\py{Field} objects). \\
\bottomrule
\end{tabular}
\end{center}

\subsection*{Reading the output}

\begin{center}
\begin{tabular}{@{}P{3cm}p{9.5cm}@{}}
\toprule
\normalfont\textbf{Symbol} & \textbf{Meaning} \\
\midrule
\texttt{Delta(...)}   & Overall momentum-conservation factor. \\
\texttt{gamma(...)}   & $\gamma$ matrix. \\
\texttt{t(...)}       & Gauge generator tensor. \\
\texttt{f(...)}       & Structure constant. \\
\texttt{g(...)}       & Metric tensor. \\
\texttt{pcomp(p,mu)}  & Component of momentum $p$ along index $\mu$. \\
\bottomrule
\end{tabular}
\end{center}

% --------------------------------------------------------------
\section{Vertex Discovery and Filtering}
\label{sec:feynpy-discovery}
% --------------------------------------------------------------

Call \py{feynman_rule()} with no arguments to enumerate every available
vertex as a dictionary.  Use the lighter-weight inspection methods to
browse without computing all expressions first.

\begin{minted}{python}
# Enumerate all vertices -- returns dict keyed by name tuples
all_rules = L.feynman_rule(simplify=False)
for sig, expr in all_rules.items():
    print(sig, "->", expr)

# Object-keyed dictionary
all_rules = L.feynman_rule(simplify=False, key_format="fields")
\end{minted}

\subsection*{Lightweight inspection with \texttt{vertex\_signatures}}
\supervisor{What is 'arity'? it appears in the code below.}
Here the \emph{arity} of a vertex means the number of external fields, or equivalently external legs, attached to it. Thus a two-point term has arity $2$, a three-point interaction has arity $3$ and so on. 
\begin{minted}{python}
# All signatures
for sig in L.vertex_signatures():
    print(sig.names, "arity=", sig.arity, "terms=", sig.term_count)

# Filter by arity
three_point  = L.vertex_signatures(arity=3)

# Exact field-content match (order-independent)
exact        = L.vertex_signatures(signature=(Gluon, Quark.bar, Quark))

# All signatures containing a given field multiset
ghost_sigs   = L.vertex_signatures(contains_fields=(GhostG.bar, GhostG))
\end{minted}

\begin{minted}{python}
# vertex_report -- filtered signatures with count summary
report = L.vertex_report(arity=3)
print(report.matched_signatures, "of", report.total_signatures)
for sig in report.signatures:
    print(sig.names)
\end{minted}

\begin{tcolorbox}[notebox, title=Tip]
For large models, use \py{vertex_signatures} or \py{vertex_report} to
inspect available vertices before calling \py{feynman_rule}, which
actually computes the symbolic expressions.
\end{tcolorbox}

% --------------------------------------------------------------
\section{A Complete Example: Gauge-Fixed QCD}
\label{sec:feynpy-qcd-example}
% --------------------------------------------------------------

We opened this chapter with the gauge-fixed QCD Lagrangian and promised to
derive its rules. All the ingredients are now available, so here is the
complete script, from imports to a compiled Lagrangian.

\noindent To keep the example compact, we import the public \textsc{FeynPy}
surface with \py{from feynpy import *}. The non-\textsc{FeynPy} objects
(\py{S}, \py{Expression}, \py{I} and the \textsc{Spenso} tensor builders) are
still imported explicitly.

\begin{minted}{python}
from fractions import Fraction

from symbolica import S, Expression
from symbolic.vertex_engine import I
from symbolic.spenso_structures import gauge_generator, structure_constant
from feynpy import *

mu, nu = S("mu"), S("nu")
a, xi, gs = S("a"), S("xi"), S("gs")
quarter = Expression.num(1) / Expression.num(4)

# Fields
Quark = Field(
    "q",
    spin=Fraction(1, 2),
    self_conjugate=False,
    symbol=S("q"),
    conjugate_symbol=S("qbar"),
    indices=(SPINOR_INDEX, COLOR_FUND_INDEX),
)
Gluon = Field(
    "G",
    spin=1,
    self_conjugate=True,
    symbol=S("G"),
    indices=(LORENTZ_INDEX, COLOR_ADJ_INDEX),
)
GhostG = GhostField(
    "ghG",
    ghost_of=Gluon,
    self_conjugate=False,
    symbol=S("ghG"),
    conjugate_symbol=S("ghGbar"),
    indices=(COLOR_ADJ_INDEX,),
)

# Gauge group
FUND = GaugeRepresentation(
    index=COLOR_FUND_INDEX,
    generator_builder=gauge_generator,
    name="fundamental",
)
SU3C = GaugeGroup(
    name="SU3C",
    abelian=False,
    coupling=gs,
    gauge_boson=Gluon,
    ghost_field=GhostG,
    structure_constant=structure_constant,
    representations=(FUND,),
)

# Gauge-fixed QCD Lagrangian
qcd = Model(
    (
        I * Quark.bar * Gamma(mu) * DC(Quark, mu)
        - quarter * FS(SU3C, mu, nu, a) * FS(SU3C, mu, nu, a)
        + GaugeFixing(SU3C, xi=xi)
        + GhostLagrangian(SU3C)
    ),
    gauge_groups=(SU3C,),
)
L = qcd.lagrangian()
\end{minted}

\noindent Before printing large expressions, inspect what the compiler
generated:

\begin{minted}{python}
print("compiled terms:", len(L.terms))

for sig in L.vertex_signatures():
    names = str(sig.names)
    print(f"{names:<25} arity={sig.arity}  terms={sig.term_count}")
\end{minted}

\noindent The output is:

\begin{minted}{text}
compiled terms: 14

('G', 'G')                arity=2  terms=5
('ghG.bar', 'ghG')        arity=2  terms=1
('q.bar', 'q')            arity=2  terms=1
('G', 'G', 'G')           arity=3  terms=4
('ghG.bar', 'G', 'ghG')   arity=3  terms=1
('q.bar', 'q', 'G')       arity=3  terms=1
('G', 'G', 'G', 'G')      arity=4  terms=1
\end{minted}

\noindent This listing already tells a small story. A single \py{FS} factor
squared has produced two-, three- and four-gluon signatures at once; the
\py{GaugeFixing} declaration has added four further terms to the two-gluon
signature; and the \py{GhostLagrangian} has contributed both a ghost kinetic
term and a ghost--gluon vertex.

\noindent Now ask for the individual rules. Since simplification is the
default, it does not need to be passed explicitly:

\begin{minted}{python}
V_qqg = L.feynman_rule(Quark.bar, Quark, Gluon)
V_cg  = L.feynman_rule(GhostG.bar, GhostG, Gluon)
V_ggg = L.feynman_rule(Gluon, Gluon, Gluon)
V_gg  = L.feynman_rule(Gluon, Gluon)
\end{minted}

\subsection*{Selected output}

\noindent\textbf{Quark--gluon.}
\begin{minted}{text}
I*gs*gamma(bis(4,i1),bis(4,i2),mink(4,mu3))
   *t(coad(8,a3),cof(3,c1),cof(3,c2))
\end{minted}
\noindent This is $i g_s \gamma^{\mu_3}t^{a_3}$, with the colour indices in
the fundamental slots of the two quark legs.

\noindent\textbf{Ghost--gluon.}
\begin{minted}{text}
gs*f(coad(8,a1),coad(8,a3),coad(8,a2))*pcomp(q1,mu3)
\end{minted}
\noindent This is $g_s f^{a_1a_3a_2}q_1^{\mu_3}$, carrying the momentum of the
antighost leg as expected from $\mathcal{L}_{\mathrm{gh}}$.

\noindent\textbf{Three-gluon.}
\begin{minted}{text}
-gs*g(mink(4,mu3),mink(4,mu1))*f(coad(8,a1),coad(8,a2),coad(8,a3))
    *pcomp(q1,mu2)
+gs*g(mink(4,mu3),mink(4,mu1))*f(coad(8,a1),coad(8,a2),coad(8,a3))
    *pcomp(q3,mu2)
+gs*g(mink(4,mu3),mink(4,mu2))*f(coad(8,a1),coad(8,a2),coad(8,a3))
    *pcomp(q2,mu1)
-gs*g(mink(4,mu3),mink(4,mu2))*f(coad(8,a1),coad(8,a2),coad(8,a3))
    *pcomp(q3,mu1)
+gs*g(mink(4,mu1),mink(4,mu2))*f(coad(8,a1),coad(8,a2),coad(8,a3))
    *pcomp(q1,mu3)
-gs*g(mink(4,mu1),mink(4,mu2))*f(coad(8,a1),coad(8,a2),coad(8,a3))
    *pcomp(q2,mu3)
\end{minted}
\noindent This is the textbook expression
\begin{equation}
    g_s f^{a_1a_2a_3}\Big[
      g^{\mu_1\mu_2}(q_1-q_2)^{\mu_3}
    + g^{\mu_2\mu_3}(q_2-q_3)^{\mu_1}
    + g^{\mu_3\mu_1}(q_3-q_1)^{\mu_2}\Big],
\end{equation}
with all momenta incoming, written out term by term.

\noindent\textbf{Gluon two-point.}
\begin{minted}{text}
 I*g(coad(8,a1),coad(8,a2))*pcomp(q1,mu1)*pcomp(q2,mu2)/xi
+I*g(mink(4,mu1),mink(4,mu2))*g(coad(8,a1),coad(8,a2))
    *pcomp(q1,mu1_int)*pcomp(q2,mu1_int)
-I*g(coad(8,a1),coad(8,a2))*pcomp(q1,mu2)*pcomp(q2,mu1)
\end{minted}
\noindent This rule shows the gauge parameter explicitly. Setting $\xi=1$
collapses the inverse propagator to the Feynman-gauge form. The label
\py{mu1_int} is an automatically generated dummy Lorentz
index, the contraction $q_1\!\cdot\!q_2$, and its name carries no
meaning, which is exactly the point of the canonicalisation discussed in
\Cref{sec:internals-postprocessing}.

% --------------------------------------------------------------
\section{Flavor Expansion}
\label{sec:feynpy-flavor-expand}
% --------------------------------------------------------------

Pass \py{flavor_expand=True} to resolve all flavor indices and produce
one Feynman rule per generation combination.  Alternatively, pass a
single \py{IndexType} or a tuple to expand only selected indices.

\begin{minted}{python}
# Compact (generation-generic) rules
rules_compact  = L.feynman_rule()

# Fully expanded -- one rule per member combination
rules_expanded = L.feynman_rule(flavor_expand=True)

# Expand only one specific index
rules_gen      = L.feynman_rule(flavor_expand=Generation_1)

# Direct member vertex (requires flavor_expand)
vertex_e_mu    = L.feynman_rule(e.bar, mu, Phi, flavor_expand=True)

# Signature discovery also supports flavor_expand
sigs = L.vertex_signatures(flavor_expand=True)
\end{minted}

\begin{tcolorbox}[warnbox, title={Off-diagonal zeroes}]
If a parameter is declared with \py{components={(i,j): 0}} for
$i \neq j$, off-diagonal member combinations are automatically absent
from the expanded rules, no manual filtering is needed.
\end{tcolorbox}

% --------------------------------------------------------------
\section{Field Transformations: From the Gauge Basis to the Physical Basis}
\label{sec:feynpy-transformations}
% --------------------------------------------------------------

Let's see how to apply field transformations. Let's refresh the field transformation for the standard model. The Lagrangian is written in terms of $B_\mu$, $W^I_\mu$ and the doublet $\Phi$, but the vertices we want involve the photon, the $Z$, the $W^\pm$, the Higgs and the Goldstone bosons.


The bridge is a set of substitutions, and \textsc{FeynPy} implements it with \py{FieldTransformation}.

\subsection*{The basic pattern}

\begin{minted}{python}
from feynpy import FieldTransformation

# B_mu  ->  -sw * Z_mu + cw * A_mu
neutral_mixing = FieldTransformation(B, -sw * Z + cw * A)

L_physical = model.transform_fields(neutral_mixing, repeat=False)
\end{minted}

\noindent Applied to a scalar current $ (D_\mu H)^\dagger (D^\mu H)$, this
moves the hypercharge current onto the physical fields. Before the
transformation the only nonzero rule is
\[
    \Gamma(H^\dagger, H, B) =
      -\tfrac{i}{2} g_1 (q_1 - q_2)^{\mu_3},
\]
and afterwards it has been replaced by
\[
    \Gamma(H^\dagger, H, A) = -\tfrac{i}{2} c_W g_1 (q_1 - q_2)^{\mu_3},
    \qquad
    \Gamma(H^\dagger, H, Z) = +\tfrac{i}{2} s_W g_1 (q_1 - q_2)^{\mu_3},
\]
while $\Gamma(H^\dagger, H, B)$ no longer exists, because $B$ is no longer a
field of the model.

\begin{tcolorbox}[infobox, title={Compile first, transform second}]
Transformations act on an \emph{already compiled} Lagrangian. Covariant
derivatives and field strengths are expanded in the original gauge basis,
where the gauge group data is meaningful, and only then are the fields
substituted. Writing \py{model.transform_fields(...)} performs both steps in
that order. This is the same ordering that \textsc{FeynRules} uses, and it is
the reason \py{DC} is never itself transformed.
\end{tcolorbox}

\subsection*{Component rules}

Component rules are the important extra ingredient for the Standard Model.
The dictionary \py{components={slot: value}} says that the rule only applies
to one fixed component of an indexed field. For example, the weak gauge field
has slots \py{(Lorentz, WeakAdj)}, so \py{components={1: 3}} selects
$W^3_\mu$.

\begin{minted}{python}
from models.SM import standard_model_weak_tensor_components

component_lagrangian = model.lagrangian().expand_index_components(
    WEAK_FUND_INDEX,
    WEAK_ADJ_INDEX,
    tensor_components=standard_model_weak_tensor_components(),
)

transformations = (
    # neutral gauge-boson mixing
    FieldTransformation(B, -sw * Z + cw * A),
    FieldTransformation(Wi, INV_SQRT2 * W.bar + INV_SQRT2 * W,
                        components={1: 1}),
    FieldTransformation(Wi, INV_SQRT2 / I * W.bar - INV_SQRT2 / I * W,
                        components={1: 2}),
    FieldTransformation(Wi, cw * Z + sw * A,
                        components={1: 3}),

    # Higgs doublet components
    FieldTransformation(Phi, -I * GP, components={0: 1}),
    FieldTransformation(Phi, vev * INV_SQRT2 + INV_SQRT2 * H
                             + I * INV_SQRT2 * G0,
                        components={0: 2}),
)

L_physical = component_lagrangian.transform_fields(
    *transformations,
    repeat=False,
)
\end{minted}

\noindent The keyword \py{repeat=False} makes the rules act as one
simultaneous \textsc{FeynRules}-style definitions block: a field created by
one rule is not transformed again by another rule in the same pass. This is
what the electroweak rotation needs.

\subsection*{Projectors and flavour rotations}

The implementation chiral projectors and flavour rotations follows the same idea. These matrix factors are attached to the appropriate spinor and flavour indices automatically:

\begin{minted}{python}
from feynpy import ProjM, ProjP, rotation

ckm = rotation(CKM, CKMDag)

fermion_transformations = (
    FieldTransformation(LL, ProjM * vl,       components={1: 1}),
    FieldTransformation(LL, ProjM * l,        components={1: 2}),
    FieldTransformation(lR, ProjP * l),
    FieldTransformation(QL, ProjM * uq,       components={1: 1}),
    FieldTransformation(QL, ckm * ProjM * dq, components={1: 2}),
    FieldTransformation(uR, ProjP * uq),
    FieldTransformation(dR, ProjP * dq),
)
\end{minted}

\noindent The complete Standard Model implementation, including the corresponding ghost mixings and the exact
ordering used in validation, can be seen in 
\path{models/SM/SM.py} by the interested user. 

% --------------------------------------------------------------
\section{Symbolica Export}
\label{sec:feynpy-symbolica-export}
% --------------------------------------------------------------

Now that we learned how to extract Feynman rules, with and without a flavor expansion, we are interested in procceding to the manipulation of the actual lagrangian. FeynPy is build on this Model API to handle all the information needed to expand the covariant derivative, field strength and so on that depends on the symmetry group. To manipulate a constructed Lagrangian it would be smart to have the expression of Symbolica so that we can call all the built-in function of this amazing library.

Call \py{to_symbolica()} on a compiled \py{Lagrangian} to obtain an \py{Expression} that can be differentiated, collected, or passed to any Symbolica routine.

\begin{minted}{python}
expr = L.to_symbolica()
print(expr)

# Standard Symbolica operations work directly
expr.expand()
expr.derivative(S("Phi"))
expr.coefficient(S("lam"))
\end{minted}

\subsection*{Coordinate-dependent export}

\begin{minted}{python}
# Makes spacetime dependence explicit
expr_coord = L.to_symbolica(derivative_style="coordinate")
expr_coord.derivative(S("x_mu"))
\end{minted}

\subsection*{Flavor-expanded export}

\begin{minted}{python}
compact  = L_fl.to_symbolica()              # generation-generic
expanded = L_fl.to_symbolica(flavor_expand=True)  # per member (e, mu, ta, ...)
\end{minted}

% --------------------------------------------------------------
\section{Applying Operators}
\label{sec:feynpy-apply-operator}
% --------------------------------------------------------------

The manipulation we can do on our \py{Model} goes further. We are interested in applying linear and non-linear operators to every term of our Lagrangian. Operators such as a replace operator, a derivative, or an infinitesimal Gauge Transformation. We will see that we can even apply a non-linear operator like the infinitesimal BRST transformation.

The operator layer acts on a compiled \py{Lagrangian} term by term,
following a graded Leibniz rule. It performs field redefinitions,
substitutions, gauge variations, and BRST transformations.

\subsection*{Required imports}

\begin{minted}{python}
from lagrangian.operator_action import (
    FieldOperator, TermOperator,
    OperatorAtomResult, OperatorSummand,
    OperatorExpansionError,
    replacement_operator, single_field_result,
    partial,
)
\end{minted}

\subsection*{Simple field replacement}

\begin{minted}{python}
op    = replacement_operator("Phi_to_Chi", {Phi: Chi()})
L_new = L.apply_operator(op)

# Works on Model too -- apply before compiling
L_new = model.apply_operator(op)
\end{minted}

\subsection*{FieldOperator --- slot-wise action}

Acts once per matched field slot, expanding via the graded Leibniz rule. As an example it's exposed the operator $O[\phi] = a*\chi + b*\Omega$

\begin{minted}{python}
# Phi -> a*Chi + b*Omega  (one summand per slot)
op = FieldOperator(
    "phi_to_combo",
    on_field=lambda occ: (
        None if occ.field is not Phi else
        OperatorAtomResult(summands=(
            OperatorSummand(coefficient=S("a"), replacement=(Chi(),)),
            OperatorSummand(coefficient=S("b"), replacement=(Omega(),)),
        ))
    ),
)

# Odd operator (parity=1 triggers graded Leibniz signs)
s = FieldOperator("s", parity=1, on_field=lambda occ: ...)
\end{minted}

\subsection*{TermOperator --- whole-term action}
We can apply an operator to an entire lagrangian term.
\begin{minted}{python}
# Acts on each InteractionTerm as a whole (no Leibniz expansion)
from dataclasses import replace
scale2 = TermOperator(
    "scale2",
    apply_to_term=lambda term: (replace(term, coupling=2 * term.coupling),),
)
\end{minted}

\subsection*{Spacetime derivative operator \texorpdfstring{--- \texttt{partial(mu)}}{}}
As a linear operator one of the most interesting is the derivative operator.
\begin{minted}{python}
L_partial = L.apply_operator(partial(mu))          # partial_mu on every slot
L_partial = L.apply_operator(partial(mu, on=Phi))  # restrict to Phi
\end{minted}

\subsection*{Composition and algebra}

\begin{minted}{python}
# Apply a sequence left-to-right
L_out = L.apply_operators(op1, op2, op3)

# Graded commutator [A, B] = A otime B - (-1)^{|A||B|} B otime A
L.operator_bracket(opA, opB).to_symbolica()
\end{minted}

\begin{tcolorbox}[infobox, title={Derivative distribution}]
When a replacement maps a field to a product, any derivative acting on that
field is automatically distributed over the product according to the Leibniz
rule. No manual handling is required.
\end{tcolorbox}

% --------------------------------------------------------------
\section{Integration by Parts}
\label{sec:feynpy-ibp}
% --------------------------------------------------------------

Two Lagrangians which differ by a total derivative give rise to the same physics.
Deciding whether that is the case is exactly the question that arises when
reducing a redundant operator basis to a minimal one
(\Cref{sec:theory-green-basis}), and \py{ibp_normal_form()} answers a
restricted version of it: it rewrites a compiled Lagrangian into a canonical
representative modulo total derivatives, so that two Lagrangians with the
Normal forms that are equivalent differ only by boundary terms.

\begin{minted}{python}
# Check that  phi^2 partial partial phi  ~  -2 phi (d_mu phi)^2  modulo total derivatives
e1 = c * phi * phi * PartialD(PartialD(phi, mu), mu)
e2 = -2 * c * phi * PartialD(phi, mu) * PartialD(phi, mu)

assert Model(e1).lagrangian().ibp_normal_form().terms \
    == Model(e2).lagrangian().ibp_normal_form().terms

# Equivalently, the normal form of the difference is empty
delta = Model(e1 - e2).lagrangian().ibp_normal_form()
assert delta.terms == ()
\end{minted}

\begin{tcolorbox}[warnbox, title={Scope of the current implementation}]
This first version of the integration-by-parts layer handles unlabelled scalar species only: it symmetrises identical scalar slots and moves derivatives into a left-normal form. Terms containing Dirac bilinears or gauge indices are outside its scope, so it cannot yet be used to reduce a Green basis to the Warsaw basis. \Cref{sec:conclusion-future} returns to
this.
\end{tcolorbox}


% --------------------------------------------------------------
\section{Gauge Variation}
\label{sec:feynpy-gauge-variation}
% --------------------------------------------------------------

We are now interested in gauge transformations. We can apply an infinitesimal gauge transformation to check the gauge invariance of a given Lagrangian. \py{gauge_variation(group=..., parameter=...)} returns a \py{FieldOperator} implementing the infinitesimal gauge variation for a givern group. Different possible transformation are exposed in the following table.


\subsection*{Transformation rules}

\begin{center}
\begin{tabular}{@{}p{5cm}p{8cm}@{}}
\toprule
\textbf{Field} & \textbf{Variation $\delta$} \\
\midrule
Matter rep $R$: $\Phi_i$           & $+ig\,\alpha^a T^a_{ij}\,\Phi_j$ \\
Matter rep $R$: $\bar\Phi_i$        & $-ig\,\alpha^a \bar\Phi_j T^a_{ji}$ \\
U(1) charge $q$: $\Phi$            & $+igq\,\alpha\,\Phi$ \\
Abelian gauge boson $A_\mu$         & $+\partial_\mu\alpha$ \\
Non-abelian gauge boson $A^a_\mu$   & $+\partial_\mu\alpha^a - g f^{abc}\alpha^b A^c_\mu$ \\
\bottomrule
\end{tabular}
\end{center}

\begin{minted}{python}
from lagrangian.operator_action import gauge_variation
from symbolic.tensor_canonicalization import canonize_full

delta = gauge_variation(group=U1QED, parameter="alpha")
dL    = L.apply_operator(delta)

# Verify invariance
canon = canonize_full(dL.to_symbolica().expand(), lorentz_indices=(mu,))
assert canon.to_canonical_string() == "0"   # => gauge invariant
\end{minted}

% --------------------------------------------------------------
\section{BRST Transformations}
\label{sec:feynpy-brst}
% --------------------------------------------------------------

The natural next step after the infinitesimal gauge transformation is to consider the BRST transformation. The constructor

\begin{minted}{python}
brst_transformation(group, ghost, antighost=None, auxiliary=...)
\end{minted}

\noindent returns an \emph{odd} \py{FieldOperator} implementing the BRST
transformation for one selected gauge group. The ghost must be declared with
\py{kind="ghost"} and associated with the chosen gauge boson through
\py{ghost_of=<gauge_boson>}.



\noindent The Brst transforamtion are 
\begin{align*}
    s A_\mu^a &= D_\mu c^a = \partial_\mu c^a + g f^{abc} A_\mu^b c^c,\\
    s c^a &= -\frac{g}{2} f^{abc} c^b c^c,\\
    s \bar c^a &= B^a,\\
    s B^a &= 0.
\end{align*}

\noindent For matter fields, we use
$$
s \Phi_i = + i g\, c^a (T^a)_{ij} \Phi_j,
\qquad
s \bar\Phi_i = - i g\, c^a \bar\Phi_j (T^a)_{ji},
$$
and, in the abelian case,
$$
s \Phi = + i g q\, c\, \Phi,
\qquad
s \bar\Phi = - i g q\, c\, \bar\Phi.
$$

If a field carries charges under several groups, we apply one BRST operator per group and sum the resulting contributions.

Because the operator is Grassmann-odd, it acts on products through the graded Leibniz rule.
Keep in mind that the exported Symbolica expression is still commutative, therefore there will be a possible sign error in how the symbolica expression is printed out.
\begin{minted}{python}
from lagrangian.operator_action import brst_transformation

# Ghost and auxiliary field declarations
c    = Field("c", spin=0, kind="ghost", ghost_of=Gluon,
             self_conjugate=False, conjugate_symbol=S("cbar"),
             indices=(COLOR_ADJ_INDEX,),
             quantum_numbers={"GhostNumber": 1})
Baux = Field("B", spin=0, self_conjugate=True, indices=(COLOR_ADJ_INDEX,))

s = brst_transformation(group=SU3C, ghost=c, auxiliary=Baux)

# Act on individual fields
mu, a = S("mu"), S("a")
print(Gluon(mu, a).apply_operator(s).to_symbolica())  # s A^a_mu
print(c(a).apply_operator(s).to_symbolica())          # s c^a
\end{minted}

\subsection*{Nilpotency check $s^2 = 0$}
The BRST transformation has a peculiar property, namly the nilpotency. We can check this property bu applying twice the operator. For example on the Gluon Field.
\begin{minted}{python}
s2_G  = Gluon(mu, a).apply_operator(s).apply_operator(s).to_symbolica()
canon = canonize_full(s2_G, run_gamma=False, run_color=False, infer_indices=True)
assert canon.to_canonical_string() == "0"
\end{minted}

\subsection*{BRST-exact gauge fixing}

We can write a BRST Lagrangian as a BRST variation of some funciton $K_{\mathrm{BRST}}$

$$
\mathcal{L}_{\mathrm{BRST}}= s(K_{\mathrm{BRST}}), \quad K_{\mathrm{BRST}}
=
\bar c^a\,\partial^\mu A^a_\mu
+
\frac{\xi}{2}\,\bar c^a B^a .
$$

Its BRST variation gives the gauge-fixing and ghost terms:
$$
sK_{\mathrm{BRST}}
=
B^a\,\partial^\mu A^a_\mu
+
\frac{\xi}{2}\,B^a B^a
-
\bar c^a\,\partial^\mu D_\mu c^a ,
$$
with
$$
D_\mu c^a
=
\partial_\mu c^a + g f^{abc} A_\mu^b c^c .
$$



The full gauge-fixed BRST-invariant Lagrangian is then
$$
L_{\mathrm{gf}}
=
L_{\mathrm{YM}} + sK_{\mathrm{BRST}} .
$$


\begin{minted}{python}
xi = S("xi")
K = Model(
    gauge_groups=(SU3C,),
    lagrangian_decl=(
        c.bar(a) * PartialD(Gluon(mu, a), mu)
        + xi / Expression.num(2) * c.bar(a) * Baux(a)
    ),
).lagrangian()

L_gf = L_YM + K.apply_operator(s)   # full gauge-fixed Lagrangian
\end{minted}

\begin{tcolorbox}[warnbox, title={Ghost field validation}]
\py{brst_transformation} validates that the ghost is declared with
\py{kind="ghost"} and that \py{ghost_of} matches the group's gauge boson.
A mismatch raises \py{ValueError} at construction time.
\end{tcolorbox}

% --------------------------------------------------------------
\section{Expression Manipulation}
\label{sec:feynpy-expr-manip}
% --------------------------------------------------------------

Once exported via \py{to_symbolica()}, standard Symbolica methods
(\py{expand}, \py{collect}, \py{coefficient}) work directly.

\begin{minted}{python}
expr = L.to_symbolica()
expr.expand()
expr.coefficient(S("xi"))      # works for literal symbols
expr.coefficient(S("gS"))      # coupling-constant dependence
\end{minted}

\subsection*{Wildcard-aware coefficient extraction}

\begin{minted}{python}
from lagrangian.symbolica_export import pattern_coefficient

mu1_, a1_, mu2_, a2_ = S("mu1_", "a1_", "mu2_", "a2_")  # trailing _ = wildcard
pattern       = G(mu1_, a1_) * G(mu2_, a2_)

print(pattern_coefficient(expr, pattern))  # handles wildcards automatically
\end{minted}



\begin{tcolorbox}[infobox, title={Wildcard convention}]
In Symbolica, variables whose names end with \texttt{\_} are interpreted as
wildcards in pattern-matching operations such as \py{match(...)}. The
resulting matches can be substituted into a pattern using
\py{replace_wildcards(...)}. This convention is useful for extracting
index-independent patterns from exported tensor expressions.
\end{tcolorbox}

% --------------------------------------------------------------
\section{Packaged Models}
\label{sec:feynpy-models}
% --------------------------------------------------------------

Up to this point, only general-purpose tools have been described. The toolkit also comes with three complete models, which are built entirely from these tools. These models are used as examples in the validation section (see \Cref{ch:validation}). Each model has its own directory under \path{models/}, along with its \textsc{FeynRules} reference data, comparison adapter, and tests.

\subsection*{The Standard Model}

\begin{minted}{python}
from models.SM import build_standard_model

sm = build_standard_model()
L  = sm.lagrangian

wwa = L.feynman_rule(sm.fields.W.bar, sm.fields.W, sm.fields.A)
hww = L.feynman_rule(sm.fields.H,     sm.fields.W.bar, sm.fields.W)
\end{minted}

\noindent The builder follows the pipeline of this chapter: declare the gauge-basis fields ($Q_L$, $L_L$, $\Phi$, $B$, $W^I$, $G$ and the singlets), assemble the gauge, fermion, Higgs, Yukawa and ghost sectors, compile the covariant derivatives and field strengths, expand the weak indices into components, and finally apply the field transformations of \Cref{sec:feynpy-transformations} in a single simultaneous pass. It returns a \py{StandardModel} record exposing the source model, the transformation list, the compiled physical-basis Lagrangian and the field and parameter namespaces.


\subsection*{The unbroken background-field Standard Model}

\begin{minted}{python}
from models.UnbrokenSM_BFM import build_unbroken_sm_bfm

bfm = build_unbroken_sm_bfm()
\end{minted}

\noindent This model stays in the unbroken gauge basis and instead performs
the background/quantum split of \Cref{sec:theory-bfm}, again as a
\py{FieldTransformation}. Because the gauge is fixed only for the quantum
fields, and covariantly with respect to the background ones, its gauge-fixing
term is written out explicitly rather than through the generic
\py{GaugeFixing} helper.

\subsection*{Green-basis SMEFT}

\begin{minted}{python}
from models.SMEFT import build_smeft_green_bpreserving

smeft = build_smeft_green_bpreserving()

Ltot  = smeft.lagrangians["Ltot"]    # EFT contribution only (the model default)
Lfull = smeft.lagrangians["Lfull"]   # Standard Model plus EFT

L = smeft.model.lagrangian()         # compiles Ltot
\end{minted}

\noindent This is the largest model, and the one that motivated most of the engine's generality. It declares $298$ Wilson coefficients as indexed \py{Parameter}s and $32$ Lagrangian blocks in the baryon-number-preserving Green basis of \Cref{sec:theory-green-basis}, including dual field strengths, $\sigma^{\mu\nu}$ structures, weak $\varepsilon$ tensors and the explicit charge-conjugation matrices of the evanescent sector. The Standard Model core is taken from the unbroken background-field model above, minus its background-field ghost and gauge-fixing sectors. Two conventions are worth noting because \Cref{ch:validation} depends on them. \texttt{Ltot} contains the EFT contribution alone, matching the \textsc{FeynRules} convention for the reference export, with the combined theory available separately as \texttt{Lfull}. And the model is written in terms of chiral fields (\py{lL}, \py{qL}, \py{eR}, \py{uR}, \py{dR}) so chirality is carried by the field names rather than by explicit projectors.

% --------------------------------------------------------------
\section{Quick Reference}
\label{sec:feynpy-quickref}
% --------------------------------------------------------------

\noindent
Here \texttt{L} denotes a compiled \py{CompiledLagrangian}, typically from
\texttt{model.lagrangian()}.

\subsection*{Construction}

\begin{center}
\begin{tabular}{@{}P{5.4cm}p{7cm}@{}}
\toprule
\normalfont\textbf{Call} & \textbf{Purpose} \\
\midrule
\texttt{IndexType(name, rep, kind)} & Define a custom index type. \\
\texttt{flavor\_index(name, dim, ...)} & Define a flavor index type. \\
\texttt{Field(name, spin, self\_conjugate, ...)} & Declare a particle field. \\
\texttt{GhostField(name, ghost\_of=..., ...)} & Declare a ghost field. \\
\texttt{GaugeRepresentation(index, generator\_builder, ...)} & Define a matter representation. \\
\texttt{GaugeGroup(name, abelian, coupling, ...)} & Define a gauge group. \\
\texttt{Parameter(name, symbol=None, indices=(), ...)} & Declare a parameter or indexed coupling tensor. \\
\bottomrule
\end{tabular}
\end{center}

\subsection*{Model and compilation}

\begin{center}
\begin{tabular}{@{}P{5.4cm}p{7cm}@{}}
\toprule
\normalfont\textbf{Call} & \textbf{Purpose} \\
\midrule
\texttt{Model(expr, ...)} & Shorthand model declaration from source terms. \\
\texttt{model.lagrangian()} & Compile; returns \py{CompiledLagrangian}. \\
\texttt{model.validate()} & Source-level validation report. \\
\bottomrule
\end{tabular}
\end{center}

\subsection*{Declarative factors (inside \texttt{lagrangian\_decl})}

\begin{center}
\begin{tabular}{@{}P{5.4cm}p{7cm}@{}}
\toprule
\normalfont\textbf{Call} & \textbf{Purpose} \\
\midrule
\texttt{Gamma(mu)} & Declarative Dirac $\gamma^\mu$ placeholder. \\
\texttt{PartialD(field, mu)} & Ordinary $\partial_\mu\,\text{field}$. \\
\texttt{DC(field, mu)} & Covariant derivative with convention $D_\mu=\partial_\mu-igA_\mu$ (alias \texttt{CovD}). \\
\texttt{FS(group, mu, nu[,a])} & $F_{\mu\nu}$; adjoint index required for non-abelian groups (alias \texttt{FieldStrength}). \\
\texttt{GaugeFixing(group, xi=...)} & Linear covariant gauge-fixing declaration. \\
\texttt{GhostLagrangian(group)} & Ordinary non-abelian Faddeev--Popov ghost declaration. \\
\bottomrule
\end{tabular}
\end{center}

\subsection*{Extraction}

\begin{center}
\begin{tabular}{@{}P{5.4cm}p{7cm}@{}}
\toprule
\normalfont\textbf{Call} & \textbf{Purpose} \\
\midrule
\texttt{L.feynman\_rule(F1, F2, ...)} & One vertex for those external fields. \\
\texttt{L.feynman\_rule()} & All discovered vertices as a dict. \\
\texttt{L.vertex\_signatures(...)} & Grouped signatures without vertex computation. \\
\texttt{L.vertex\_report(...)} & Structured filtered report with grouped counts. \\
\bottomrule
\end{tabular}
\end{center}

\subsection*{Compiled Lagrangian layer}

\begin{center}
\begin{tabular}{@{}P{5.4cm}p{7cm}@{}}
\toprule
\normalfont\textbf{Call} & \textbf{Purpose} \\
\midrule
\texttt{L.to\_symbolica(...)} & Export / display as a Symbolica \py{Expression}. \\
\texttt{L.transform\_fields(*rules)} & Substitute fields, e.g.\ gauge basis $\to$ physical basis. \\
\texttt{L.expand\_index\_components(...)} & Expand finite gauge indices into explicit components. \\
\texttt{L.apply\_operator(op)} & Apply a \py{FieldOperator} or \py{TermOperator}. \\
\texttt{L.apply\_operators(...)} & Compose operators left-to-right. \\
\texttt{L.operator\_bracket(A, B)} & Graded commutator $[A,B]$. \\
\texttt{L.ibp\_normal\_form()} & Scalar-only IBP normal form. \\
\bottomrule
\end{tabular}
\end{center}

\subsection*{Operator constructors}

\begin{center}
\begin{tabular}{@{}P{5.4cm}p{7cm}@{}}
\toprule
\normalfont\textbf{Call} & \textbf{Purpose} \\
\midrule
\texttt{replacement\_operator(name, \{F: r\})} & Simple field-replacement operator. \\
\texttt{FieldOperator(name, parity, ...)} & Slot-wise operator with graded Leibniz rule. \\
\texttt{TermOperator(name, ...)} & Whole-term operator. \\
\texttt{partial(mu)} & Runtime spacetime-derivative operator. \\
\texttt{gauge\_variation(group, parameter=...)} & Infinitesimal gauge variation $\delta_\alpha$. \\
\texttt{brst\_transformation(group, ghost, ...)} & Odd BRST differential. \\
\texttt{FieldTransformation(source, expr, ...)} & Declarative field redefinition. \\
\texttt{rotation(V, Vdag)} & Flavour-rotation matrix factor for a replacement. \\
\bottomrule
\end{tabular}
\end{center}

\subsection*{Packaged models}

\begin{center}
\begin{tabular}{@{}P{5.4cm}p{7cm}@{}}
\toprule
\normalfont\textbf{Call} & \textbf{Purpose} \\
\midrule
\texttt{build\_standard\_model()} & Standard Model in the physical basis. \\
\texttt{build\_unbroken\_sm\_bfm()} & Unbroken Standard Model with the background-field split. \\
\texttt{build\_smeft\_green\_bpreserving()} & Green-basis SMEFT, EFT-only by default. \\
\bottomrule
\end{tabular}
\end{center}

\subsection*{Expression manipulation}

\begin{center}
\begin{tabular}{@{}P{5.4cm}p{7cm}@{}}
\toprule
\normalfont\textbf{Call} & \textbf{Purpose} \\
\midrule
\texttt{expr.expand()} / \texttt{.coefficient(s)} & Standard Symbolica algebra. \\
\texttt{pattern\_coefficient(expr, pat)} & Wildcard-aware coefficient extraction. \\
\texttt{canonize\_full(expr, ...)} & Full tensor simplification and canonicalisation. \\
\bottomrule
\end{tabular}
\end{center}
